\section{Protein}

\subsection{Khái niệm và đặc điểm cấu tạo}

Protein là những hợp chất hữu cơ có phân tử khối rất lớn, được tạo thành từ nhiều đơn vị amino acid liên kết với nhau. Trong phân tử protein có liên kết peptide \(\ce{-CO-NH-}\).

Có thể mô tả sự hình thành liên kết peptide qua phản ứng ngưng tụ giữa hai amino acid:
\[
\ce{H2N-CHR1-COOH + H2N-CHR2-COOH
-> H2N-CHR1-CO-NH-CHR2-COOH + H2O}.
\]

Phân tử protein thường có cấu trúc phức tạp; các amino acid khác nhau và trật tự sắp xếp của chúng tạo nên tính chất riêng của từng protein.

\subsection{Tính chất hóa học}

\textbf{Phản ứng thủy phân}

Khi đun nóng protein trong môi trường acid hoặc base, hoặc dưới tác dụng của enzyme thích hợp, các liên kết peptide bị cắt, tạo hỗn hợp amino acid.

\[
\text{Protein} + \ce{H2O}
\xrightarrow[\text{đun nóng}]{\text{acid/base/enzyme}}
\text{hỗn hợp amino acid}.
\]

\textbf{Sự đông tụ}

Protein có thể bị đông tụ khi tác dụng với acid, base hoặc khi đun nóng. Ví dụ: lòng trắng trứng chuyển từ dạng lỏng sang dạng đông đặc khi đun nóng.

Khi đun nóng mạnh, protein có thể bị phân hủy và tạo ra các sản phẩm có mùi đặc trưng.

\subsection{Phân biệt một số loại sợi}

Len lông cừu và tơ tằm có bản chất protein; tơ nylon là polymer tổng hợp, không phải protein.

Có thể khảo sát một lượng rất nhỏ sợi bằng phép thử đốt trong điều kiện an toàn:
\begin{itemize}
\item Sợi len hoặc tơ tằm có xu hướng cháy sém, co lại và tạo mùi giống tóc/lông cháy.
\item Sợi nylon thường nóng chảy, co thành hạt cứng và có mùi nhựa cháy.
\end{itemize}
Chỉ thực hiện với mẫu rất nhỏ, dùng kẹp, tránh hít trực tiếp khói và tuân thủ hướng dẫn an toàn của giáo viên.

\subsection{Vai trò của protein}

Protein là thành phần quan trọng cấu tạo nên tế bào và mô cơ thể; tham gia tạo enzyme, một số hormone và nhiều chất cần thiết cho hoạt động sống. Cơ thể cần được cung cấp protein từ nhiều nguồn thực phẩm phù hợp.

\begin{ghinho}{Liên hệ thực tế}
Thịt, cá, trứng, sữa, đậu và một số loại hạt là những nguồn thực phẩm giàu protein. Cần phối hợp đa dạng thực phẩm để cung cấp đủ chất dinh dưỡng cho cơ thể.
\end{ghinho}
